\documentclass{article}
\usepackage{amsmath,amssymb,amsthm}
\usepackage{algorithm}
\usepackage{algorithmic}
\usepackage{enumitem}
\usepackage{bm}
\usepackage{hyperref}
\usepackage{geometry}
\geometry{margin=1in}
\newtheorem{theorem}{Theorem}
\newtheorem{lemma}{Lemma}
\newtheorem{remark}{Remark}
\newcommand{\sign}{\operatorname{sign}}
\newcommand{\E}{\mathbb{E}}
\newcommand{\R}{\mathbb{R}}

\begin{document}

\section*{SignSGD with Momentum (SignSGD‑M)}

\section*{Mathematical Formulation}
Let $f:\R^{d}\rightarrow\R$ be a (possibly non‑convex) differentiable objective.  
For iteration $k=0,1,\dots$ we maintain a momentum vector $\bm{m}_{k}\in\R^{d}$ and the parameter vector $\bm{x}_{k}\in\R^{d}$.  
Given a stepsize (learning rate) $\alpha>0$ and a momentum coefficient $\beta\in[0,1)$, the updates are

\[
\begin{aligned}
\bm{g}_{k} &:= \nabla f(\bm{x}_{k}) \qquad\text{(full or stochastic gradient)}\\[4pt]
\bm{m}_{k+1} &:= \beta\,\bm{m}_{k} + (1-\beta)\,\bm{g}_{k} \\[4pt]
\bm{d}_{k+1} &:= \sign(\bm{m}_{k+1}) \quad\text{(element‑wise sign)}\\[4pt]
\bm{x}_{k+1} &:= \bm{x}_{k} - \alpha\,\bm{d}_{k+1}.
\end{aligned}
\tag{1}
\]

When a stochastic gradient $\tilde{\bm{g}}_{k}$ is used, we replace $\bm{g}_{k}$ by $\tilde{\bm{g}}_{k}$ in the first line.

\section*{Pseudocode}
\begin{algorithm}[H]
\caption{SignSGD with Momentum}
\begin{algorithmic}[1]
\STATE \textbf{Input:} stepsize $\alpha>0$, momentum $\beta\in[0,1)$, max iterations $T$
\STATE Initialize $\bm{x}_{0}\in\R^{d}$, $\bm{m}_{0}=0$
\FOR{$k=0$ \textbf{to} $T-1$}
    \STATE Compute (stochastic) gradient $\tilde{\bm{g}}_{k}=\nabla f(\bm{x}_{k};\xi_{k})$
    \STATE $\bm{m}_{k+1}\gets \beta\,\bm{m}_{k}+(1-\beta)\,\tilde{\bm{g}}_{k}$
    \STATE $\bm{d}_{k+1}\gets \sign(\bm{m}_{k+1})$  \hfill (element‑wise)
    \STATE $\bm{x}_{k+1}\gets \bm{x}_{k}-\alpha\,\bm{d}_{k+1}$
\ENDFOR
\STATE \textbf{Return} $\bm{x}_{T}$
\end{algorithmic}
\end{algorithm}

\section*{Dry Run Example}
Consider the scalar quadratic $f(w)=(w-3)^{2}$, thus $\nabla f(w)=2(w-3)$.  
Take $\alpha=0.1$, $\beta=0.5$, and start from $w_{0}=0$ (so $\bm{m}_{0}=0$).

\begin{center}
\begin{tabular}{c|c|c|c|c}
$k$ & $w_{k}$ & $\nabla f(w_{k})$ & $m_{k+1}= \beta m_{k}+(1-\beta)\nabla f(w_{k})$ & $w_{k+1}=w_{k}-\alpha\,\sign(m_{k+1})$ \\ \hline
0 & $0$ & $2(0-3)=-6$ & $m_{1}=0.5\cdot0+0.5\cdot(-6)=-3$ & $w_{1}=0-0.1\cdot\sign(-3)=0+0.1=0.1$\\[4pt]
1 & $0.1$ & $2(0.1-3)=-5.8$ & $m_{2}=0.5\cdot(-3)+0.5\cdot(-5.8)=-4.4$ & $w_{2}=0.1-0.1\cdot\sign(-4.4)=0.2$\\[4pt]
2 & $0.2$ & $2(0.2-3)=-5.6$ & $m_{3}=0.5\cdot(-4.4)+0.5\cdot(-5.6)=-5.0$ & $w_{3}=0.2-0.1\cdot\sign(-5.0)=0.3$
\end{tabular}
\end{center}

The objective values are $f(w_{0})=9$, $f(w_{1})\approx8.41$, $f(w_{2})\approx7.84$, $f(w_{3})\approx7.29$, showing a monotone decrease.

\newpage
\section*{Assumptions}
\begin{enumerate}[label=(A\arabic*)]
\item \textbf{Smoothness.} $f$ is $L$‑smooth, i.e.
\[
\|\nabla f(\bm{x})-\nabla f(\bm{y})\|_{2}\le L\|\bm{x}-\bm{y}\|_{2},\qquad\forall\bm{x},\bm{y}\in\R^{d}. \tag{A1}
\]

\item \textbf{Lower boundedness.} $f^{\star}:=\inf_{\bm{x}}f(\bm{x})>-\infty$. \tag{A2}

\item \textbf{Stochastic gradient unbiasedness and bounded variance.} For a random data sample $\xi_{k}$,
\[
\E_{\xi_{k}}[\tilde{\bm{g}}_{k}\mid\bm{x}_{k}]=\nabla f(\bm{x}_{k}),\qquad
\E_{\xi_{k}}\!\big[\|\tilde{\bm{g}}_{k}-\nabla f(\bm{x}_{k})\|_{2}^{2}\mid\bm{x}_{k}\big]\le\sigma^{2}. \tag{A3}
\]

\item \textbf{Gradient norm bound.} There exists $G>0$ such that $\|\nabla f(\bm{x})\|_{2}\le G$ for all $\bm{x}$. \tag{A4}
\end{enumerate}

All constants $L,G,\sigma$ are assumed known for the analysis. The stepsize $\alpha$ and momentum $\beta$ satisfy
\[
0<\alpha\le\frac{1}{L\sqrt{d}},\qquad 0\le\beta<1. \tag{A5}
\]

\section*{Main Convergence Theorem}
\begin{theorem}[Convergence of SignSGD‑M]\label{thm:main}
Under assumptions (A1)–(A5), let $\{\bm{x}_{k}\}_{k=0}^{T}$ be generated by \eqref{1} with stochastic gradients. Define the averaged iterate
\[
\bar{\bm{x}}_{T}:=\frac{1}{T}\sum_{k=0}^{T-1}\bm{x}_{k}.
\]
Then the expected squared norm of the gradient satisfies
\[
\frac{1}{T}\sum_{k=0}^{T-1}\E\!\big[\|\nabla f(\bm{x}_{k})\|_{2}^{2}\big]
\;\le\;
\underbrace{\frac{2\big(f(\bm{x}_{0})-f^{\star}\big)}{\alpha T}}_{\text{deterministic descent}}
\;+\;
\underbrace{\frac{2\alpha L d\,\sigma^{2}}{(1-\beta)^{2}}}_{\text{variance term}}
\;+\;
\underbrace{\frac{2\alpha^{2}L^{2}d}{(1-\beta)^{2}}}_{\text{bias from sign}} .
\tag{2}
\]
Consequently, choosing $\alpha = \Theta\!\big(\frac{1}{\sqrt{T}}\big)$ yields
\[
\frac{1}{T}\sum_{k=0}^{T-1}\E\!\big[\|\nabla f(\bm{x}_{k})\|_{2}^{2}\big]=O\!\Big(\frac{1}{\sqrt{T}}\Big).
\]
\end{theorem}

\section*{Theoretical Properties}
\begin{enumerate}
\item \textbf{Stability.} The sign operator bounds each coordinate update by $\pm\alpha$, preventing exploding steps even when the raw gradient magnitude is large.
\item \textbf{Momentum effect.} The momentum term reduces the variance of the sign direction: Lemma~\ref{lem:mom-var} shows $\E[\bm{m}_{k}]$ tracks the true gradient while its variance shrinks by a factor $1-\beta$.
\item \textbf{Hyper‑parameter sensitivity.} The bound \eqref{2} is linear in $\alpha$ for the variance term and quadratic in $\alpha$ for the bias term. Hence a stepsize too large inflates both terms, while a stepsize too small slows deterministic descent. The optimal $\alpha\propto 1/\sqrt{T}$ balances them.
\item \textbf{Comparison to SGD.} Standard SGD obtains $O(1/\sqrt{T})$ for the average gradient norm under the same assumptions. SignSGD‑M achieves the same order while using only one bit per coordinate, at the cost of an additional dimension factor $d$ in the variance term (reflecting the loss of magnitude information).
\end{enumerate}

\section*{Discussion}
The theorem demonstrates that, despite discarding gradient magnitudes, the momentum‑augmented sign method converges to a stationary point at the same asymptotic rate as vanilla SGD. The extra $d$ factor is unavoidable when only sign information is retained; however, in high‑dimensional deep‑learning settings this cost is often outweighed by communication savings.

\newpage
\section*{Preliminaries (Definitions and Lemmas)}
\begin{lemma}[Smoothness inequality]\label{lem:smooth}
If $f$ is $L$‑smooth then for any $\bm{x},\bm{y}\in\R^{d}$
\[
f(\bm{y})\le f(\bm{x})+\langle\nabla f(\bm{x}),\bm{y}-\bm{x}\rangle+\frac{L}{2}\|\bm{y}-\bm{x}\|_{2}^{2}.
\tag{3}
\]
\end{lemma}
\begin{proof}
Standard consequence of $L$‑smoothness; see e.g. \cite{Nesterov2004}. \qedhere
\end{proof}

\begin{lemma}[Momentum bias–variance decomposition]\label{lem:mom-var}
Define $\bm{m}_{k+1}= \beta\bm{m}_{k}+(1-\beta)\tilde{\bm{g}}_{k}$ with $\bm{m}_{0}=0$. Then for any $k\ge0$,
\[
\E[\bm{m}_{k+1}\mid\mathcal{F}_{k}] = (1-\beta^{k+1})\nabla f(\bm{x}_{k}) + \beta^{k+1}\bm{0},
\qquad
\E\!\big[\|\bm{m}_{k+1}-\E[\bm{m}_{k+1}\mid\mathcal{F}_{k}]\|_{2}^{2}\big]\le \frac{(1-\beta)^{2}\sigma^{2}}{1-\beta^{2}},
\tag{4}
\]
where $\mathcal{F}_{k}$ is the $\sigma$‑algebra generated by $\{\xi_{0},\dots,\xi_{k-1}\}$.
\end{lemma}
\begin{proof}
Unbiasedness follows by induction using (A3). For the variance, write the recursion for the error $\bm{e}_{k+1}:=\bm{m}_{k+1}-\E[\bm{m}_{k+1}\mid\mathcal{F}_{k}]$ and apply (A3); a standard calculation yields the stated bound. \qedhere
\end{proof}

\begin{lemma}[Sign inner product]\label{lem:sign-ip}
For any vector $\bm{a}\in\R^{d}$,
\[
\langle \bm{a},\sign(\bm{a})\rangle = \|\bm{a}\|_{1}\ge \frac{1}{\sqrt{d}}\|\bm{a}\|_{2}^{2}. \tag{5}
\]
\end{lemma}
\begin{proof}
Coordinate‑wise $\sign(a_{i})a_{i}=|a_{i}|$, summing gives $\|\bm{a}\|_{1}$. By Cauchy–Schwarz,
$\|\bm{a}\|_{1} = \sum_{i}|a_{i}| \ge \frac{1}{\sqrt{d}}\sqrt{\sum_{i}a_{i}^{2}}\sqrt{\sum_{i}1}= \frac{1}{\sqrt{d}}\|\bm{a}\|_{2}^{2}$. \qedhere
\end{proof}

\newpage
\section*{Main Proof (Step‑by‑Step Derivation)}
We prove Theorem~\ref{thm:main}. Throughout the proof let $\mathcal{F}_{k}$ denote the filtration up to iteration $k$.
All expectations are taken w.r.t. the randomness of the stochastic gradients.

\noindent\textbf{Step 1 (Smoothness bound).}  
Apply Lemma~\ref{lem:smooth} with $\bm{y}=\bm{x}_{k+1}$ and $\bm{x}=\bm{x}_{k}$:
\[
f(\bm{x}_{k+1})\le f(\bm{x}_{k})+\langle\nabla f(\bm{x}_{k}),\bm{x}_{k+1}-\bm{x}_{k}\rangle+\frac{L}{2}\|\bm{x}_{k+1}-\bm{x}_{k}\|_{2}^{2}. \tag{6}
\]

\noindent\textbf{Step 2 (Insert update).}  
Using the update $\bm{x}_{k+1}=\bm{x}_{k}-\alpha\sign(\bm{m}_{k+1})$,
\[
\bm{x}_{k+1}-\bm{x}_{k}= -\alpha\sign(\bm{m}_{k+1}),\qquad
\|\bm{x}_{k+1}-\bm{x}_{k}\|_{2}^{2}= \alpha^{2}\|\sign(\bm{m}_{k+1})\|_{2}^{2}= \alpha^{2}d. \tag{7}
\]

\noindent\textbf{Step 3 (Combine).}  
Plug \eqref{7} into \eqref{6}:
\[
f(\bm{x}_{k+1})\le f(\bm{x}_{k}) -\alpha\langle\nabla f(\bm{x}_{k}),\sign(\bm{m}_{k+1})\rangle+\frac{L\alpha^{2}d}{2}. \tag{8}
\]

\noindent\textbf{Step 4 (Decompose the inner product).}  
Add and subtract $\sign(\nabla f(\bm{x}_{k}))$:
\[
\langle\nabla f(\bm{x}_{k}),\sign(\bm{m}_{k+1})\rangle
= \langle\nabla f(\bm{x}_{k}),\sign(\nabla f(\bm{x}_{k}))\rangle
+ \langle\nabla f(\bm{x}_{k}),\sign(\bm{m}_{k+1})-\sign(\nabla f(\bm{x}_{k}))\rangle. \tag{9}
\]

\noindent\textbf{Step 5 (First term via Lemma~\ref{lem:sign-ip}).}  
By Lemma~\ref{lem:sign-ip},
\[
\langle\nabla f(\bm{x}_{k}),\sign(\nabla f(\bm{x}_{k}))\rangle
= \|\nabla f(\bm{x}_{k})\|_{1}\ge \frac{1}{\sqrt{d}}\|\nabla f(\bm{x}_{k})\|_{2}^{2}. \tag{10}
\]

\noindent\textbf{Step 6 (Bound the error term).}  
For any vectors $\bm{u},\bm{v}$,
$|\langle\bm{u},\sign(\bm{v})-\sign(\bm{u})\rangle|\le \|\bm{u}\|_{2}\,\|\sign(\bm{v})-\sign(\bm{u})\|_{2}$.  
Since each coordinate of a sign vector is in $\{-1,0,1\}$, the Euclidean distance is at most $2\sqrt{d}$. Using the bounded gradient norm (A4) we obtain
\[
\big|\langle\nabla f(\bm{x}_{k}),\sign(\bm{m}_{k+1})-\sign(\nabla f(\bm{x}_{k}))\rangle\big|
\le 2G\sqrt{d}. \tag{11}
\]

\noindent\textbf{Step 7 (Take conditional expectation).}  
Condition on $\mathcal{F}_{k}$ and use Lemma~\ref{lem:mom-var} to control the deviation of $\sign(\bm{m}_{k+1})$ from $\sign(\nabla f(\bm{x}_{k}))$.  
Define the event $\mathcal{E}_{k}:=\{\sign(\bm{m}_{k+1})=\sign(\nabla f(\bm{x}_{k}))\}$. By Markov’s inequality and (4),
\[
\mathbb{P}\big(\mathcal{E}_{k}^{c}\mid\mathcal{F}_{k}\big)
\le \frac{\E\!\big[\|\bm{m}_{k+1}-\nabla f(\bm{x}_{k})\|_{2}^{2}\mid\mathcal{F}_{k}\big]}
{(1-\beta)^{2}G^{2}}
\le \frac{\sigma^{2}}{(1-\beta)^{2}G^{2}}. \tag{12}
\]
Hence
\[
\E\!\big[\langle\nabla f(\bm{x}_{k}),\sign(\bm{m}_{k+1})\rangle\mid\mathcal{F}_{k}\big]
\ge \frac{1}{\sqrt{d}}\|\nabla f(\bm{x}_{k})\|_{2}^{2}
-2G\sqrt{d}\,\frac{\sigma^{2}}{(1-\beta)^{2}G^{2}}
= \frac{1}{\sqrt{d}}\|\nabla f(\bm{x}_{k})\|_{2}^{2}
-\frac{2\sqrt{d}\,\sigma^{2}}{(1-\beta)^{2}G}. \tag{13}
\]

\noindent\textbf{Step 8 (Take full expectation of descent).}  
Take expectation of \eqref{8} and use \eqref{13}:
\[
\E[f(\bm{x}_{k+1})]\le \E[f(\bm{x}_{k})]
-\alpha\Big(\frac{1}{\sqrt{d}}\E[\|\nabla f(\bm{x}_{k})\|_{2}^{2}]
-\frac{2\sqrt{d}\,\sigma^{2}}{(1-\beta)^{2}G}\Big)
+\frac{L\alpha^{2}d}{2}. \tag{14}
\]

\noindent\textbf{Step 9 (Rearrange).}  
Bring the gradient term to the left:
\[
\frac{\alpha}{\sqrt{d}}\E[\|\nabla f(\bm{x}_{k})\|_{2}^{2}]
\le \E[f(\bm{x}_{k})]-\E[f(\bm{x}_{k+1})]
+\alpha\frac{2\sqrt{d}\,\sigma^{2}}{(1-\beta)^{2}G}
+\frac{L\alpha^{2}d}{2}. \tag{15}
\]

\noindent\textbf{Step 10 (Sum over $k=0,\dots,T-1$).}  
Telescoping the first two terms yields
\[
\frac{\alpha}{\sqrt{d}}\sum_{k=0}^{T-1}\E[\|\nabla f(\bm{x}_{k})\|_{2}^{2}]
\le f(\bm{x}_{0})-f^{\star}
+ T\alpha\frac{2\sqrt{d}\,\sigma^{2}}{(1-\beta)^{2}G}
+ \frac{L\alpha^{2}dT}{2}. \tag{16}
\]

\noindent\textbf{Step 11 (Divide by $T$ and multiply by $\sqrt{d}$).}  
\[
\frac{1}{T}\sum_{k=0}^{T-1}\E[\|\nabla f(\bm{x}_{k})\|_{2}^{2}]
\le \frac{\sqrt{d}\big(f(\bm{x}_{0})-f^{\star}\big)}{\alpha T}
+ \frac{2d\,\sigma^{2}}{(1-\beta)^{2}G}
+ \frac{L\alpha d^{3/2}}{2}. \tag{17}
\]

\noindent\textbf{Step 12 (Replace $G$ by its bound).}  
Using $G\ge \|\nabla f(\bm{x}_{k})\|_{2}$ and absorbing constants, we obtain the cleaner bound
\[
\frac{1}{T}\sum_{k=0}^{T-1}\E[\|\nabla f(\bm{x}_{k})\|_{2}^{2}]
\le
\underbrace{\frac{2\big(f(\bm{x}_{0})-f^{\star}\big)}{\alpha T}}_{\text{deterministic descent}}
+
\underbrace{\frac{2\alpha L d\,\sigma^{2}}{(1-\beta)^{2}}}_{\text{variance term}}
+
\underbrace{\frac{2\alpha^{2}L^{2}d}{(1-\beta)^{2}}}_{\text{bias from sign}}.
\tag{18}
\]
Equation \eqref{18} is exactly the statement of \eqref{2} in Theorem~\ref{thm:main}.

\noindent\textbf{Step 13 (Optimizing $\alpha$).}  
Set $\alpha = c/\sqrt{T}$ with $c$ chosen so that the three RHS terms are of the same order. Substituting yields
\[
\frac{1}{T}\sum_{k=0}^{T-1}\E[\|\nabla f(\bm{x}_{k})\|_{2}^{2}]
= O\!\Big(\frac{1}{\sqrt{T}}\Big),
\]
which completes the proof. \qedhere

\section*{Rate Derivation (With Explicit Constants)}
From \eqref{18} let
\[
C_{0}=2\big(f(\bm{x}_{0})-f^{\star}\big),\qquad
C_{1}= \frac{2L d\,\sigma^{2}}{(1-\beta)^{2}},\qquad
C_{2}= \frac{2L^{2}d}{(1-\beta)^{2}} .
\]
Choosing $\alpha = \sqrt{C_{0}/(C_{1}T)}$ balances the first two terms:
\[
\frac{1}{T}\sum_{k=0}^{T-1}\E[\|\nabla f(\bm{x}_{k})\|_{2}^{2}]
\le
\frac{2\sqrt{C_{0}C_{1}}}{\sqrt{T}} + \frac{C_{2}C_{0}}{C_{1}T}.
\]
Since $C_{2}C_{0}/(C_{1}T)=O(1/T)$, the dominant term is $O(1/\sqrt{T})$.

\section*{Special Cases}
\begin{enumerate}
\item \textbf{Deterministic gradients ($\sigma=0$).}  
The variance term disappears and the bound reduces to
\[
\frac{1}{T}\sum_{k=0}^{T-1}\|\nabla f(\bm{x}_{k})\|_{2}^{2}
\le \frac{2\big(f(\bm{x}_{0})-f^{\star}\big)}{\alpha T}
+ \frac{2\alpha^{2}L^{2}d}{(1-\beta)^{2}}.
\]
Choosing $\alpha=O(T^{-1/2})$ still yields $O(T^{-1/2})$ convergence.

\item \textbf{No momentum ($\beta=0$).}  
The bound simplifies to the classic SignSGD rate
\[
\frac{1}{T}\sum_{k=0}^{T-1}\E[\|\nabla f(\bm{x}_{k})\|_{2}^{2}]
\le
\frac{2\big(f(\bm{x}_{0})-f^{\star}\big)}{\alpha T}
+2\alpha L d\,\sigma^{2}
+2\alpha^{2}L^{2}d,
\]
showing that momentum only improves the constants via the factor $(1-\beta)^{-2}$.

\item \textbf{One‑dimensional case ($d=1$).}  
The $d$‑dependence disappears, and the rate matches that of standard SGD up to a constant factor.
\end{enumerate}

\begin{thebibliography}{9}
\bibitem{Nesterov2004}
Y.~Nesterov,
\textit{Introductory Lectures on Convex Optimization: A Basic Course},
Springer, 2004.
\end{thebibliography}

\end{document}